%\frontmatter
\thispagestyle{empty}
\section*{Resumen}

Las integrinas son proteínas integrales de membranas. Son grandes receptores de superficie celular heterodiméricos ($\alpha$\!$\beta$) que desempeñan un papel clave en la formación de complejos de adhesión a la matriz extracelular y están involucrados en varias vías de transducción de señales. La integrina $\alpha$IIb$\beta$3 está presente en plaquetas y es el principal modulador del proceso de agregación de plaquetas. En el estado desactivado, los segmentos transmembrana de la integrina $\alpha$IIb y de $\beta$3 forman un heterodímero con orientación relativa definida. La disociación del heterodímero y la consecuente activación del receptor esta mediada por la unión de ligandos citoplasmáticos a segmento intracelular de $\beta$3. Las interacciones que establecen los segmentos transmembrana en el heterodímero son determinantes para la funcionalidad del receptor. En está tesis empleamos monocapas de Langmuir como sistema de membrana modelo, y simulaciones de dinámica molecular en bicapas lipídicas para estudiar la función del ambiente lipídico en la conformación y orientación de los segmentos transmembrana e intracelular de la integrina $\alpha$IIb$\beta$3.

De estos estudios surgieron una serie de ideas novedosas, entre ellas: (i) los péptidos presentan una orientación vertical respecto al plano de las monocapas de Langmuir; (ii) el análisis de polaridad sugiere que el segmento intracelular está en contacto con la subfase acuosa y el segmento transmembrana queda expuesto hacia el aire; (iii) la presencia de los péptidos modifica las propiedades reológicas de la membrana como la tensión superficial, elasticidad lateral y módulo de compresibilidad; (iv) se observó que el segmento de $\alpha$IIb interacciona preferentemente con lípidos en estado líquido expandido (POPC o la mezcla POPC-POPG) que un lípido en estado líquido condensado (DPPC); (v) en general, el ambiente lipídico modula las interacciones y la orientación entre los segmentos transmembrana del heterodímero. 

Los resultados experimentales y bioinformáticos presentados aquí demuestran la importancia de comprender las propiedades del entorno en el que se encuentran las proteínas de membrana y el impacto que la naturaleza del ambiente lipídico tiene sobre sus conformaciones y funcionalidad.


\newpage
\thispagestyle{empty}
\section*{Abstract}

Integrins are examples of integral membrane proteins. They are large heterodimeric cell surface receptors ($\alpha$-$\beta$) that play a key role in forming adhesion complexes with the extracellular matrix and are involved in several signal transduction pathways. Integrin $\alpha$IIb$\beta$3 is found in platelets and is the primary modulator of platelet aggregation. In its inactive state, the transmembrane segments of integrin $\alpha$IIb and $\beta$3 form a heterodimer with a specific relative orientation. The dissociation of this heterodimer, leading to receptor activation, is triggered by the binding of cytoplasmic ligands to the intracellular segment of $\beta$3. The interactions established by the transmembrane segments in the heterodimer are crucial for the functionality of the receptor. In this thesis, we use Langmuir monolayers as a model membrane system, along with molecular dynamics simulations in lipid bilayers, to study the role of the lipid environment in the conformation and orientation of the transmembrane and intracellular segments of the integrin $\alpha$IIb$\beta$3.

From these studies, a series of novel ideas emerged, including: (i) the peptides have a vertical orientation with respect to the plane of the Langmuir monolayers; (ii) Polarity analysis suggests the intracellular segment is in contact with the aqueous subphase and the transmembrane segment is exposed to the air; (iii) the presence of the peptides alters the rheological properties of the membrane, such as surface tension, lateral elasticity, and compressibility modulus; (iv) it was observed that lipids such as POPC and PCPG preferentially interact with the $\alpha$IIb peptide; and (v) in general, the lipid environment modulates the interactions and orientation between the transmembrane segments of the heterodimer..

The experimental and bioinformatics results presented here demonstrate the importance of understanding the properties of the environment in which membrane proteins are found, as well as the impact that the nature of the lipid environment has on their conformations and functionality.
